\section{Deferral and Boundary}

The decision for X.3 is to defer both RLOG and RCHAIN. No RLOG crate, RCHAIN
crate, production type, parser, fixture ladder, or verifier command is claimed
here. The downstream-boundary and civic-evidence access-pattern specifications
both place these packages after the base evidence layers and state that RLOG and
RCHAIN should be added only when real log and custody source artifacts demand
them \citep{civic-evidence-boundaries-2026}.

That order matters. Public EMS logs, scanner logs, tabulator logs, upload logs,
adjudication logs, seal logs, custody forms, ballot retrieval sheets, and
observer records vary by jurisdiction and vendor. A schema written before those
artifacts are collected would mostly encode guesses. X.3 therefore requires L0
source-pressure fixtures before any crate work: small positive and negative
examples based on real public record shapes, with source bytes or source-like
records, expected normalizations, and explicit caveats.

The boundary rule is the central contribution: RLOG can say what a source log
records. RCHAIN can say what a custody source records. Neither can infer
custody or correctness from count reconciliation alone. If an RCOUNT package
balances perfectly, that is evidence about the count contract. It is not by
itself evidence that a seal was intact, that a room was continuously observed,
that every machine event was exported, or that no custody gap occurred.

The same rule also limits future composition. A later RCASE bundle may cite
RCOUNT, RLOG, and RCHAIN together, and RCERT may reference them when an official
action depends on process evidence. But composition must preserve the child
claim boundaries. It may say which package supports which public claim; it must
not let a count verifier silently become a custody verifier.
